\documentclass[11pt]{article}
\usepackage[a4paper,margin=24mm]{geometry}
\usepackage{amsmath,amssymb,amsthm,booktabs,longtable,array}
\usepackage[T1]{fontenc}
\usepackage[colorlinks=true,urlcolor=blue]{hyperref}
\newtheorem{theorem}{Theorem}
\newtheorem{lemma}[theorem]{Lemma}
\newtheorem{proposition}[theorem]{Proposition}
\newcommand{\CE}{C_{\mathrm E}}
\newcommand{\Q}{\mathbb Q}
\title{A zero-loss redesign of the upper-index proof\\
\large Conditional composition and a small exact support certificate}
\author{Research audit in the fork \texttt{kolodziejek/berry-esseen-certificates}}
\date{27 October 2026}
\begin{document}
\maketitle
\begin{abstract}
A first failure after a sharp finite prefix has zero predecessor decrement.
A support-free scalar recurrence then supplies the local score gates without
positive-ceiling propagation. We derive a new support certificate, uniform
for all actual indices $n\ge42$, using 14 closed intervals and 672 rational
Bernstein coefficients. This removes the finite upper endpoint from the
support part of the proposed LOCAL42 route. The resulting strict-tail theorem
is conditional on a genuine sharp base and the retained analytic and local
numerical implications. It is not an acceptance of the full conjecture.
\end{abstract}

\section{Scope, notation and retained assumptions}
For standardized iid summands with finite third absolute moment
$\rho=E|X|^3$, write
\[
 C_n=\sup_P\frac{\sqrt n}{\rho(P)}
       \sup_x\left|P\{(X_1+\cdots+X_n)/\sqrt n\le x\}-\Phi(x)\right|,
 \qquad \CE=\frac{3+\sqrt{10}}{6\sqrt{2\pi}}<\frac{41}{100}.
\]
The inequality for $\CE$ and the treatment of reflected left limits are
recorded in the independent scalar audit. Attainment of an oriented positive
closed-CDF maximizing pair when $C_n\ge\CE$ is a separate retained analytic
input; this includes equality, not merely strict exceedance.

The proposed uniform LOCAL42 statement is:
\begin{equation}\label{local}
\begin{gathered}
 n\ge42,\quad C_n\ge\CE,\quad
 0\le n(C_n-\CE)\le .54,\quad n(C_{n-1}-C_n)_+=0,\\
 C_n<.4226,\qquad C_{n-1},C_{n-2}<.4253
 \quad\Longrightarrow\quad\text{contradiction}.
\end{gathered}
\end{equation}
Here \eqref{local} is conditional on the complete finite 42 resource and
comparison predicates and on their analytic licenses, with the new support
certificate below replacing the old finite-index support signs.
The detailed transfer of the other index-sensitive hypotheses is in
\texttt{uniform42.md}. It retains weighted resources $nD\le42d_D$,
same-law joint gap/noise constraints, both signs, strict antichain widths,
mode floor 40, actual strips and moment coefficients, signed normalization
repairs, the complete closed TINY/MID/LOW cover, and the zero-noise Bernoulli
theorem. None is proved by the scalar recurrence or by a support calculation.

\section{The scalar entry and first-failure composition}
\begin{lemma}[Support-free smoothing]
For every integer $n\ge2$,
\[
 C_n\le \sqrt{\frac n{n-1}}C_{n-1}
 +\frac{\sqrt n}{15(n-1)^{3/2}}+\frac{2\sqrt n}{25(n-1)^2}.
\]
\end{lemma}
\begin{proof}
Put $v=n-1$ and $F_v(t)=\Phi(t/\sqrt v)$. Conditioning on the last summand
introduces at most $C_{n-1}\rho/\sqrt v$. Taylor's formula through order two
and $EX=0$, $EX^2=1$ give
\[
 \left|EF_v(t-X)-F_v(t)-\tfrac12F_v''(t)\right|
 \le \frac{\rho}{15v^{3/2}},
\]
because $\|\phi''\|_\infty=\phi(0)<2/5$.
The heat identity $\partial_v F_v=F_v''/2$ gives
\[
 \left|F_{v+1}-F_v-\tfrac12F_v''\right|
 \le \frac{1}{4\pi v^2}<\frac{2}{25v^2}.
\]
Indeed $\|\phi'''\|_\infty\le2/\pi$ follows from the Gaussian Fourier
integral, and the integrated second variance derivative has factor $1/8$.
Multiply by $\sqrt n/\rho$, use $\rho\ge1$, and take suprema. The uniform
remainder uses only the third absolute moment and holds at every threshold.
\end{proof}

\begin{lemma}[Entry at 42]\label{entry}
If $n\ge42$ and $C_{n-1}\le\CE\le C_n$, then
\[
 0\le n(C_n-\CE)<\frac{3003}{10250}<\frac3{10},
 \qquad C_n<\frac{73}{175}<.4226 .
\]
\end{lemma}
\begin{proof}
For $s=\sqrt{n/(n-1)}$, rationalizing $s-1$ in the preceding lemma yields
\[
 n(C_n-\CE)\le\frac n{n-1}
 \left(\frac{\CE}{s+1}+\frac{s}{15}
                +\frac{2s}{25\sqrt{n-1}}\right).
\]
Now $s<81/80$, $\sqrt{n-1}>6$, $s+1>2$ and $n/(n-1)\le42/41$.
Thus the displayed expression is strictly smaller than
\[
 \frac{42}{41}\left(\frac{41}{200}
             +\frac{27}{400}+\frac{27}{2000}\right)
 =\frac{3003}{10250}.
\]
The reserve below $3/10$ is $36/5125$. Also
$C_n<.41+.3/42=73/175$, whose reserve below $.4226$ is $191/35000$.
\end{proof}

\begin{theorem}[Conditional strict tail]\label{tail}
Assume $C_j\le\CE$ for $1\le j\le41$, and assume the equality-inclusive
implication \eqref{local} with all of its analytic and numerical premises.
Then $C_n<\CE$ for every $n\ge42$.
\end{theorem}
\begin{proof}
If the conclusion fails, choose the least integer $n\ge42$ with $C_n\ge\CE$.
Both predecessors are at most $\CE$; when $n=42$ these are exactly 40 and 41.
Hence their caps are strictly below $.4253$, and
$n(C_{n-1}-C_n)_+=0$. Lemma~\ref{entry} supplies stronger excess and current
score bounds than \eqref{local}. The retained equality-inclusive attainment
hypothesis and \eqref{local} give a contradiction.
Choosing the least strict exceedance instead would establish only a weak tail.
\end{proof}

\section{Uniform contact polynomials}
Let $c=819/2000<\CE$ and $N=42$. The lower inequality follows from
$\sqrt{10}>3081/500-3$, $\pi<22/7$ and
$\sqrt{44/7}<1567/625$: the resulting lower bound for $\CE$ exceeds
$c$ by $377/3134000$. Let $r,b$ be upper-bound parameters satisfying
$\rho\le r$ and $|\beta|\le b$ for the same current law.
They are not lower bounds on the actual moments. Put
\[
 \ell=\frac{6403124237432}{10^{12}},\qquad
 S=\frac{202424330939}{200000000000},\qquad
 \ell^2<41,\quad S^2>\frac{42}{41}.
\]
For $u\ge5/2$ define
\begin{align*}
 A&=(u^3-3u)/6,&J&=u^5-15u,&H_4&=u^4-6u^2+3,\\
 D_+&=J/(120\cdot41),&W&=A-\tfrac{15}{2}D_+,&B_2&=H_4^2/(256\cdot41),\\
 E&=\frac{51u^6}{14400\cdot41\ell}
      +\frac{51(3u^2-2)}{960\cdot41\ell}
      +\frac{u}{8\cdot41^2},\\
 Q&=-u^3+\tfrac32ru^2+3bu+\tfrac{.384}{42}r .
\end{align*}
The infinite-upper-index, zero-loss polynomials are
\begin{align}
 G&=\left(\tfrac A2-3D_+-\frac{H_4}{12\ell}\right)W+B_2,\label{g}\\
 P&=\tfrac25S^3(AW+B_2)+S^3EW+cQW.\label{p}
\end{align}
The retained signed fifth-order contact envelope states that positivity of
$A,H_4,J,W,G,-Q,-P$ excludes support at the corresponding $u$, for both signs
and every real raw threshold. Its source is
\texttt{paper/analytic-supplement.txt}, lines 455--640.
The adverse terms are bounded at 42; the favorable fifth-order lower bound is
zero for the unbounded index range. Thus \eqref{p} has no hidden finite 48
endpoint. The zero predecessor decrement removes the contact-loss term.

For comparison, the old finite polynomial has $D_-=J/(120\cdot47)$ and
loss $q=.19$. Direct subtraction gives
\[
 P_{\infty,0}-P_{48,.19}
 =SW\left(\frac{S^2J}{4700}-.19r\right),\qquad
 G_\infty-G_{48}=3D_-W.
\]
This gives an alternative transfer from old finite signs on suitable shorter
intervals, proved in \texttt{zero-loss-support.md}. The following new
certificate proves the infinite polynomials directly and therefore does not
need those old sign premises.

\section{Fourteen exact intervals}
Group the 64 LEFT support calls by their 13 radii. Within each group use the
componentwise maximum of the recorded preceding-resource bounds $r,b$.
The initial call is a separate row. The fixed table is:
\begin{center}\small
\begin{tabular}{rrrr}
\toprule $R$&$r$&$b$&Calls\\ \midrule
$5/2$&$603319/500000$&$32128627/10^8$&17\\
$101/40$&$475374/390625$&$8224681/25000000$&5\\
$51/20$&$61183419/50000000$&$33662533/10^8$&4\\
$103/40$&$30833983/25000000$&$8698269/25000000$&5\\
$13/5$&$30994887/25000000$&$35540061/10^8$&4\\
$21/8$&$24988347/20000000$&$36644311/10^8$&5\\
$53/20$&$125900813/10^8$&$37372893/10^8$&4\\
$107/40$&$63303431/50000000$&$1201521/3125000$&4\\
$27/10$&$63753427/50000000$&$489197/1250000$&4\\
$109/40$&$128341449/10^8$&$20069857/50000000$&4\\
$11/4$&$25857241/20000000$&$41146991/10^8$&4\\
$111/40$&$129968327/10^8$&$42077471/10^8$&3\\
$121/40$&$71458741/50000000$&$46459359/10^8$&1\\
$96/25$&$7361/4000$&$641562334877/10^{12}$&1\\
\bottomrule
\end{tabular}
\end{center}
The first 13 intervals are $[R,39/10]$; the final interval is $[96/25,21/5]$.
All are closed and use a single cell. For fixed $u>0$, $W>0$,
\[
 \partial_rP=cW(3u^2/2+.384/42)>0,\qquad
 \partial_bP=3cuW>0.
\]
Thus componentwise grouping preserves an upper bound even if the two maxima
come from different calls. The first five sign polynomials do not depend on
$r,b$. Each call is used sequentially only after its own actual-law resource
bounds are acquired.

\begin{proposition}[Exact interval certificate]\label{certificate}
For every table row, the seven polynomials
\[ A,\quad H_4,\quad J,\quad W,\quad G,\quad -Q,\quad -P \]
are strictly positive on its specified closed interval.
\end{proposition}
\begin{proof}
For $p(u)=\sum_{i=0}^d a_i u^i$ on $[l,h]$, write $u=l+(h-l)t$ and
\[
 q_j=(h-l)^j\sum_{i=j}^d a_i\binom ij l^{i-j},\qquad
 B_k=\sum_{j=0}^k q_j\frac{\binom kj}{\binom dj}.
\]
The binomial theorem proves
$p(l+(h-l)t)=\sum_{k=0}^d B_k\binom dk t^k(1-t)^{d-k}$.
The basis is nonnegative and sums to1 on $[0,1]$.
The degrees are $3,4,5,5,10,3,11$, giving 48 coefficients per row.
The fixed rational checker computes all $14\cdot48=672$ coefficients,
and every coefficient exceeds $1/4$. It also independently reconstructs
\[
 q_j=\binom dj\sum_{k=0}^j(-1)^{j-k}\binom jk B_k
\]
for every polynomial. Every coefficient is printed in the acceptance log,
so the finite arithmetic is reproducible without numerical optimization.
This proves continuous-domain positivity, including both endpoints.
\end{proof}

\section{Completing the rays}
The retained direct Gaussian derivative estimate, valid for every raw $H,t$, is
\[
 |\partial_tH_n(H,t)|\le\tfrac25(t^2/2+1/2+e_n),\qquad
 e_n=n\bigl(\sqrt{n/(n-1)}-1\bigr)-1/2<.384/n.
\]
For completeness, set $a=48/125$ and $x=1/n$. Expansion gives
\[
 (1-x)(1+x/2+ax^2)^2-1
 =x^2\{2a-3/4-(a+1/4)x+(a^2-a)x^2-a^2x^3\}.
\]
The braces decrease on $[0,1/42]$ and equal
$118707/42875000>0$ at its right endpoint. Taking square roots proves the
stated uniform scale correction.

For each group set
\[
 L(u)=\tfrac25(u^2/2+1/2+.384/42)+cQ'(u).
\]
At its splice $U$, the checker verifies all six strict inequalities
$Q,Q',Q'',L,L',L''<0$. Since $Q'''=-6$ and $L''=2/5-6c<0$,
these signs persist for all $u\ge U$.
For $\eta=\pm1$ the contact envelope $H_n(H,\eta u)+cQ(u)$ has derivative
at most $L(u)<0$. Its value at $U$ is already negative by
Proposition~\ref{certificate}. Hence both full exterior rays are excluded.
For each individual call the group caps majorize $Q$ and $Q'$.
This conclusion applies to atoms, arbitrarily small masses, and every
support point, with no limiting law or bounded-threshold substitution.

Combining the closed interval with the complete rays yields the original
strict support conclusion $|X|<R$ at every actual $n\ge42$, conditional on
the retained contact/heat analysis and the actual preceding moment bounds.

\section{Dependency savings and unresolved work}
The new upper-index composition needs 19 of the old 31 finite native actions
and all 152 LOW row producers. Six positive-ceiling/support actions and six
special 43 bootstrap actions are bypassed by Lemma~\ref{entry} and sharp
predecessors. If the full uniform LOCAL42 transfer is accepted, the separate
55-action tail package is also unnecessary in this new upper-index module.
Thus its conservative count is $238\to171$, before new composition joins.
This is a task count, not measured runtime or a repository-wide saving.

The grouped support certificate replaces only the support polynomial
subchecks; it does not remove the remaining resource, confinement, profile,
TINY/MID/LOW or harmonic obligations. A separate source audit found repeated
LOW evaluations by producers and full binder replay, but no one-pass receipt
implementation is included. All 152 producers remain in the conservative plan.

A genuine BASE41 remains an input. The current mathematical floor 40 route
states a sharp 38/39-to40/41/42 induction without the old upper-index tail.
The complete pinned current provider body was unavailable in the inspected
sources, while an older global 40 inventory explicitly consumes tail 43 and
positive ceilings. Consequently neither executable base closure nor
repository-wide obsolescence of those computations is established.
Attainment at equality, published moment inequalities, analytic
contact/Fourier/normalization licenses and strict Bernoulli comparison also
remain separate premises.

\section{Reproduction and preserved reference}
The fixed checker is
\texttt{mathematics/simplification/2026-10-27/check\_grouped\_support.py}.
It imports only a new source-binding helper, never archived programs.
The first acceptance returned 0 in about 0.12 seconds, with nine rejection
controls and all exact coefficients stored under
\texttt{runs/grouped-support-20261027/accept01/}.
Other bounded checks, independent derivations, source locations and limitations
are in the adjacent Markdown reports and \texttt{reviews/} directory.

The immutable original baseline is commit
\texttt{6346fc1b4432bc5982822aa696eb566ed9e1e1ee}.
The backend archive SHA256 begins \texttt{43b606138f826b47e}; full pins are in
the reports and receipts. No inherited certificate, theorem statement,
acceptance target or production runner is changed. No native heavy
certificate program was run. This document proves new conditional
implications and a small exact support certificate, not the full conjecture.
\end{document}
